\subsection{P029 --- Operational digital twins in the built environment: A systematic review of semantic infrastructure and deployment practice}
\label{paper:P029}
\textbf{Authors:} Chady Elias, Isabella Hutchins, Raja R. A. Issa\par
\textbf{Major Category:} Reviews, Frameworks and Standards\par
\textbf{Secondary Category:} Digital Twin, BIM and Semantic Information, Retrieval, Query, LLM and Agent\par
\textbf{Keywords:} Digital Twin, Semantic interoperability, Knowledge graphs, Facility management, Systematic literature review, AI agents\par
\paragraph{主な主張}
46件のoperational Digital Twin deploymentをsystematic reviewし, formal semantic modelは25件, semanticsをruntime system logicで使用したのは12件であった. Runtimeでsemanticsを用いるdeploymentではdiagnostic outputが最も多い一方, execute-level actionは主としてsite-specificなfixed mappingに依存していた. \cite{Elias:2026OperationalDTReview}
\paragraph{新規性}
既往reviewがdefinition, application, analytics等を中心に整理してきたのに対し, semantic tier, actionability, workflow integrationを用いて実装済みoperational DTを分類し, semantic infrastructureのruntime利用とdeployment readinessをimplementation-centeredに評価した点に新規性がある.
\paragraph{位置付け}
operational DTにおけるsemantic infrastructureをportable analyticsおよびagent-enabled workflowsの前提となるlifecycle information requirementとして整理したreviewであり, scalable and AI-ready building operationの研究ギャップを示す位置付けとなる.
